\documentclass[10pt,a4paper]{amsart}
\usepackage[margin=19mm]{geometry}
\usepackage{amsmath,amssymb,mathtools}
\usepackage[scr=rsfs,cal=euler]{mathalfa}
\usepackage{xfrac}
\usepackage[unicode,hidelinks,bookmarks=false]{hyperref}
\newcommand{\ie}{i.e.\ }
\newcommand{\cf}{cf.\ }
\newcommand{\resp}{respectively\ }
\newcommand{\Subsection}{\subsection}
\providecommand{\nobreakdash}{\nobreak\hbox{-}}
\setlength{\emergencystretch}{2em}
\multlinegap0pt
\usepackage{bm}
\usepackage{bbm}
\usepackage{dsfont}
\usepackage{xspace}
\usepackage{cancel}
\usepackage{mathtools}
\usepackage{amsthm}
\makeatletter
\@ifundefined{lemma}{\newtheorem{lemma}{Lemma}}{}
\makeatother
\newcommand{\FD}{\mathrm{FD}}
\newcommand{\Ker}[0]{\mathrm{Ker}\;}
\newcommand{\R}[0]{\mathbb R}
\newcommand{\cX}{\mathcal{X}}
\newcommand{\norm}[1]{\left\lVert#1\right\rVert}
\renewcommand{\dim}[0]{\mathrm{dim}\;}
\title[What Data Enables Optimal Decisions? An Exact Characterizati]{What Data Enables Optimal Decisions? An Exact Characterization for Linear Optimization}
\author{Omar Bennouna, Amine Bennouna, Saurabh Amin, Asuman Ozdaglar}
\date{Source: arXiv:2505.21692v1 (CC BY 4.0)}
\begin{document}
\maketitle

\noindent\emph{Source and licence.} What Data Enables Optimal Decisions? An Exact Characterization for Linear Optimization. Version arXiv:2505.21692v1 (CC BY 4.0), distributed under the Creative Commons Attribution 4.0 International licence, as recorded in the arXiv licence metadata for this exact version and shown on the abstract page https://arxiv.org/abs/2505.21692v1. Full rights notice: licenses/arxiv-2505.21692-CC-BY-4.0.txt.\par\smallskip

\noindent\textit{Editorial scope.} Counted unit: the Lemma whose source label is \texttt{lemma:feasible directions in general polyhedron} in the article (source0.tex lines 1122--1124, source label \texttt{lemma:feasible directions in general polyhedron}), delivered together with the article's own complete original proof of it (source0.tex lines 1125--1150, one \texttt{proof} environment of 26 lines). No mathematics is written, repaired or re-derived: statement and proof are the article's own text line for line. Required context retained uncounted: none. Every definition, display and label of those lines is retained unchanged. Omitted from this package and not redistributed: 1--1121, 1151--1160. Nothing inside a retained excerpt is omitted or altered: every retained body is delivered exactly as the article's lines are written, so no omission receipt of any kind accompanies this package. Citations: the retained excerpts carry no citation command at all, so no bibliography entry of the article is transcribed and no reference is redistributed. Reference adaptation: none was needed and none was made; every reference inside the delivered unit resolves inside it, so no unresolved reference or citation is produced. Printed slips kept literal: no printed word, symbol, hypothesis or number of the counted statement or of its proof was altered. Layout and class: the article's own class is \texttt{article} with packages including geometry, graphicx, import, amsmath, amssymb, bbold, mathtools, footnote, acronym, enumerate, hyperref, cleveref, multirow, subcaption; the wrapper is the lane's pinned \texttt{amsart} editorial wrapper at 10pt on A4, which restates the theorem-like environments the retained text needs and adds the article's own macro definitions that the retained text needs (\texttt{\textbackslash FD}, \texttt{\textbackslash Ker}, \texttt{\textbackslash R}, \texttt{\textbackslash cX}, \texttt{\textbackslash dim}, \texttt{\textbackslash norm}), with extra wrapper packages (bm, bbm, dsfont, xspace, cancel, mathtools). The delivered numbering is the unit's own. Assets: no retained span contains a figure environment, a graphics-inclusion command, a PDF-page inclusion, a table or a TikZ picture; no image file is copied into this package, the package declares no asset dependency and no image file is redistributed with it. Licence: version arXiv:2505.21692v1 of this article is distributed under the Creative Commons Attribution 4.0 International licence, as recorded in the arXiv licence metadata for this exact version and shown on the abstract page https://arxiv.org/abs/2505.21692v1. A full rights notice is delivered at licenses/arxiv-2505.21692-CC-BY-4.0.txt. Characters: every retained excerpt is pure ASCII, so the delivered engine, XeLaTeX with its default Latin Modern text font, reports no missing character.
\begin{lemma} \label{lemma:feasible directions in general polyhedron}
            Let $r=\dim F_0 \cap \Ker A$. For any $x\in \cX$, there exists a set $V=\{\delta_1,\dots,\delta_r\}\subset \FD(x)$, such that $V$ is a basis of $F_0 \cap \Ker A$. In particular, the set of extreme directions of $\FD(x)$ spans $F_0\cap \Ker A$.
        \end{lemma}

\begin{proof}
            Let $x\in \cX$. Let $\cX^\angle$ be the set of extreme points of $\cX$, and $D^\angle$ be the set of extreme rays of $\cX$. For every $x^\angle \in \cX^\angle$ and $\delta^\angle \in D^\angle$, we have $x^\angle \geq 0$ and $\delta^\angle \geq 0$. Let $\{\alpha_{x^\angle}\}_{x^\angle \in \cX^\angle}\subset \R^*_+$ and $\{\alpha_{\delta^\angle}\}_{\delta^\angle \in D^\angle}\subset \R^*_+$ be a set of strictly positive numbers such that $\sum_{x^\angle \in \cX^\angle}^{}\alpha_{x^\angle}=1$. We define 
            \begin{align*}
                \overline{x}=\sum_{x^\angle \in \cX^\angle}^{}\alpha_{x^\angle}x^\angle + \sum_{\delta^\angle \in D^\angle}^{}\alpha_{\delta^\angle}\delta^\angle\in \cX.
            \end{align*}

        We have for any $i\in I_0$, $\overline{x}_i>0$. Indeed, for any $i\in I_0$, if $\overline{x}_i=0$, then for every $x^\angle \in \cX^\angle$ and $\delta^\angle \in D^\angle$, $x_i^\angle=\delta_i^\angle =0$ and hence for any $x'\in \cX$, $x'_i=0$ i.e. $i\not \in I_0$ which is impossible. Let $$\varepsilon:=\frac{1}{2}\min_{i\in I_0}\overline{x}_i>0,$$ for any $\delta \in F_0 \cap \Ker A$ such that $\norm{\delta}<\varepsilon$, we have $A(\overline{x}+\delta)=A\overline{x}=b$, and for every $i\in [d]$, if $i\in I_0$, then $\overline{x}_i + \delta _i >0$ and if $i\not\in I_0$, $\overline{x}_i=\delta_i=0$ and consequently $x+\delta \geq 0$, i.e. $\delta$ is a feasible direction for $\overline{x}$. Hence, every element of $B(0,\varepsilon)\cap F_0 \cap \Ker A$ is a feasible direction for $\overline{x}$, and consequently any element of $F_0 \cap \Ker A$. Let $x\in \cX$, and $v_1,\dots, v_r$ a basis of $F_0\cap \Ker A$ such that for every $i\in [r],\; \norm{v_i}=1$. Let $\eta \in \R^*_+$ small enough such that $\forall i \in [r]$, $\overline{x} + \eta \delta _i \in \cX$. We would like to show that for a well-chosen value of $\eta$, the following set of feasible directions for $x$, $\{\overline{x}+\eta v_1 - x,\dots,\overline{x}+\eta v_r - x,\}$ is a basis of $F_0 \cap \Ker A$. Since $\overline{x}-x\in F_0 \cap \Ker A$, we consider $\beta_1,\dots,\beta_r\in \R$ such that 
        \begin{align*}
            \overline{x}-x=\sum_{i=1}^{r}\beta_i v_i.
        \end{align*}
        Let $\alpha_1,\dots,\alpha_r\in \R$. We have 
        \begin{align*}
            \sum_{i=1}^{r}\alpha_i(\overline{x}+\eta v_i-x)=0 &\Longrightarrow \underbrace{\left(\sum_{i=1}^{r}\alpha_i\right)}_{:=A}(\overline{x}-x)+\eta \sum_{i=1}^{r}\alpha_i v_i=0\\
            &\Longrightarrow \sum_{i=1}^{r}(A\beta_i + \eta \alpha_i)v_i=0\\
            &\Longrightarrow \forall i \in [r],\; A\beta_i + \eta \alpha_i=0\\
            &\Longrightarrow \underbrace{ \forall i\in [r], \alpha_i=0 \text{ or }\left( A\neq 0 \text{ and } \forall i \in [r],\;\frac{\alpha_i}{A}=-\frac{\beta_i}{\eta}\right)}_{(*)}.
        \end{align*}
        By summing over $i$ in the last equality above, we get
        \begin{align*}
            (*)&\Longrightarrow  \forall i\in [r], \alpha_i=0 \text{ or }\left( A\neq 0 \text{ and } 1=-\frac{1}{\eta}\smash{\underbrace{\sum_{i=1}^{r}\beta_i}_{B}} \right)
        \end{align*}
        
        \vspace{5mm}
        
        The equality above is equivalent to $\eta = -B$. Hence, it suffices to take $\eta\neq -B$ and $\eta$ small enough to ensure that $\{\overline{x}+\eta v_1 - x,\dots,\overline{x}+\eta v_r - x,\}$ is linearly independent and is a set of feasible directions for $x$, and consequently since all of these vectors are elements of $F_0 \cap \Ker A$, and there are $r=\dim F_0 \cap \Ker A$ of them, $\{\overline{x}+\eta v_1 - x,\dots,\overline{x}+\eta v_r - x,\}$ is indeed a set of feasible directions for $x$ that is a basis of $F_0\cap \Ker A$.
        \end{proof}

\end{document}
